\begin{figure}[htbp]
\centering
\resizebox{0.86\textwidth}{!}{%
\begin{tikzpicture}[
  font=\sffamily\small,
  force/.style={-{Latex[length=2.5mm]},line width=1.0pt,draw=VIUOrange!90!black},
  aux/.style={-{Latex[length=2.2mm]},line width=0.9pt,draw=VIUDark!75},
  robot/.style={circle,draw=VIUDark!80,fill=VIUOrange!20,line width=0.9pt,minimum size=8.5mm,font=\sffamily\scriptsize\bfseries},
  note/.style={draw=VIUDark!60,fill=VIUWhite,rounded corners=4pt,align=left,text width=4.2cm,inner sep=6pt}
]
\coordinate (C) at (0,0);
\begin{scope}[rotate=14]
  \filldraw[fill=VIULight,draw=VIUDark,line width=1pt,rounded corners=2pt] (-2.65,-1.22) rectangle (2.65,1.22);
  \node[robot] (FL) at (-2.65,1.22) {FL};
  \node[robot] (FR) at (2.65,1.22) {FR};
  \node[robot] (RL) at (-2.65,-1.22) {RL};
  \node[robot] (RR) at (2.65,-1.22) {RR};
  \draw[VIUDark!45,dashed] (-2.65,0) -- (2.65,0);
  \draw[VIUDark!45,dashed] (0,-1.22) -- (0,1.22);
\end{scope}
\fill[VIUDark] (C) circle (1.8pt) node[below right=2pt,font=\scriptsize] {$q_k=(x_k,y_k,\psi_k)$};

\draw[force] (FR) -- ++(-1.25,0.42) node[midway,above=1pt] {$f_{FR}$};
\draw[force] (FL) -- ++(-0.65,0.55) node[midway,left=1pt] {$f_{FL}$};
\draw[force] (RL) -- ++(-0.72,-0.48) node[midway,left=1pt] {$f_{RL}$};
\draw[force] (RR) -- ++(0.68,-0.45) node[midway,right=1pt] {$f_{RR}$};
\draw[aux] (0.35,1.8) arc[start angle=34,end angle=142,radius=1.9] node[midway,above] {$\tau_z$};

\node[draw=VIUDark!60,fill=VIUDark!8,circle,minimum size=1.1cm] (obs) at (4.5,0.65) {obs.};
\draw[draw=VIUOrange!65,dashed,line width=0.8pt] (obs) circle (1.0cm);

\node[note] at (0,-3.45) {$W_k=\sum_\rho J_\rho(q_k)^\top f_\rho$\\[2pt]
$\tau_z$ aparece si la l\'inea de acci\'on no pasa por el centro.};
\end{tikzpicture}}
\caption[Carga rectangular con roles geométricos]{Carga rectangular con roles geométricos de coalición. Las fuerzas aplicadas en las esquinas generan un wrench sobre la carga; si no pasan por el centro, aparece un torque que permite rotar para rodear obstáculos.}
\viuownsource
\label{fig:carga-rectangular-roles}
\end{figure}
